\documentclass[11pt,reqno]{amsart}
\usepackage{amsmath,amssymb,amsthm,mathtools}
\usepackage[margin=1.15in]{geometry}
\usepackage{microtype}
\usepackage[colorlinks=true,linkcolor=blue,citecolor=blue,urlcolor=blue]{hyperref}

\theoremstyle{plain}
\newtheorem{theorem}{Theorem}[section]
\newtheorem{proposition}[theorem]{Proposition}
\newtheorem{lemma}[theorem]{Lemma}
\newtheorem{corollary}[theorem]{Corollary}
\theoremstyle{definition}
\newtheorem{remark}[theorem]{Remark}
\newtheorem{conjecture}[theorem]{Conjecture}

\newcommand{\N}{\mathbb{N}}
\newcommand{\Z}{\mathbb{Z}}
\DeclareMathOperator{\tot}{\varphi}

\begin{document}

\title[Erd\H{o}s--Graham \#411 for $r=2$]{Erd\H{o}s--Graham \#411 for $r=2$:\\
a totient bridge, a generalized-Fermat cascade,\\ and a machine-checked $\omega$-ladder}

\author{Jared Wilder}
\email{jared@jaredwilder.com}

\date{\today}

\begin{abstract}
Let $g(n)=n+\tot(n)$. Erd\H{o}s and Graham asked for the solutions of $g_{k+r}(n)=2g_k(n)$.
For $r=2$, Steinerberger reduced the problem to $\tot(n)+\tot(n+\tot(n))=n$ and showed the
solution list $n=2^{\ell}\{1,3,5,7,35,47\}$ is complete unless there is a prime
$p\equiv 7 \pmod 8$, $p>47$, with $\tot\!\big((3p-1)/4\big)=(p+1)/2$; he verified none exists
below $10^{10}$.

We give an exact bridge from that exceptional condition to the totient equation
$3\tot(N)=2N+2$ (equivalently $\tot(n)=\tfrac23(n+1)$) studied independently by Hercher, a
connection drawn in neither source. Four consequences follow.
(i) The congruence clause is vacuous: $p\equiv7\pmod 8$ holds automatically, so the entire
content of the $r=2$ case is the \emph{primality} of $p=2\tot(N)-1$.
(ii) The known solutions form a base-$6$ generalized-Fermat cascade $n_j=6^{2^j}-1$, and we
prove $n_j$ solves the equation if and only if $6^{2^k}+1$ is prime for all $k<j$; since
$6^{8}+1=17\cdot 98801$, the cascade \emph{provably terminates} at $j=3$ and contributes
exactly the primes $7$ and $47$. A local sieve, $\ell\mid p \iff N\equiv-4^{-1}\!\!\pmod \ell$,
explains precisely why the other two cascade candidates are composite.
(iii) Combining the bridge with Hercher's lower bound gives: \emph{no exceptional prime exists
below $1.33\times10^{14}$}, improving the published bound by a factor of $13{,}000$.
(iv) An $\omega$-ladder, driven by a single arity-general ``multiplicative balance'' cap
theorem proved here, gives $\omega(N)\ge 6$ machine-checked in Lean~4 and $\omega(N)\ge 8$
by certified exhaustive enumeration, past the published $\omega\ge7$.

We also prove a no-go: for every prime $\ell\ge5$ all admissible residues of $N$ are
attained, so no finite covering congruence can force $p$ composite. Congruence methods alone
cannot resolve the problem. All elementary results are formalized in Lean~4/Mathlib.
\end{abstract}

\maketitle

\section{Introduction}

Write $\tot$ for Euler's totient and $g(n)=n+\tot(n)$, with $g_k$ the $k$-fold iterate.
Problem \#411 of the Erd\H{o}s--Graham list asks for solutions of $g_{k+r}(n)=2g_k(n)$.
Steinerberger \cite{Stein} settled the structure of the $r=2$ case: it is equivalent to
\begin{equation}\label{eq:stein}
  \tot(n)+\tot\big(n+\tot(n)\big)=n ,
\end{equation}
whose solutions are $n=2^{\ell}\cdot m$ with $m\in\{1,3,5,7,35,47\}$ --- \emph{complete unless}
there exists a prime
\begin{equation}\label{eq:exceptional}
  p\equiv 7 \pmod 8,\qquad p>47,\qquad \tot\!\left(\frac{3p-1}{4}\right)=\frac{p+1}{2}.
\end{equation}
Such a $p$ we call \emph{exceptional}. Steinerberger verified that none exists with
$p\le 10^{10}$ and left the question open.

Independently and within weeks, Hercher \cite{Herch} studied the equation
$\tot(n)=\tfrac23(n+1)$, proving every solution is squarefree, chain-free, and that any
solution beyond the known set $\{5,35,1295,1679615\}$ has at least seven distinct prime
factors and exceeds $10^{14}$.

Neither paper refers to the other, and neither observes that the two problems are the same
problem. Establishing that is the starting point of this note.

\subsection*{Results}
Throughout we set
\[
  N=\frac{3p-1}{4},\qquad\text{equivalently}\qquad p=\frac{4N+1}{3}.
\]

\begin{theorem}[The bridge]\label{thm:bridge}
$p$ satisfies $\tot(N)=(p+1)/2$ if and only if $3\tot(N)=2N+2$, i.e.\ iff
$\tot(N)=\tfrac23(N+1)$. Consequently the exceptional primes of \eqref{eq:exceptional} are
exactly the primes of the form $p=(4N+1)/3$ arising from solutions $N>35$ of Hercher's
equation.
\end{theorem}

\begin{theorem}[The residue clause is free]\label{thm:free}
For \emph{any} odd $N$ with $N\equiv2\pmod3$, the integer $p=(4N+1)/3$ satisfies
$p\equiv7\pmod8$ automatically. Hence the congruence condition in \eqref{eq:exceptional}
carries no information, and the entire content of the $r=2$ case is the primality of
$p=2\tot(N)-1$.
\end{theorem}

\begin{theorem}[Cascade, and its termination]\label{thm:cascade}
Put $n_j=6^{2^j}-1$. Then $3\tot(n_j)=2n_j+2$ if and only if $6^{2^k}+1$ is prime for every
$k<j$. Since $6^{8}+1=17\cdot 98801$ is composite, the only members of the cascade solving the
equation are $n_0=5$, $n_1=35$, $n_2=1295$, $n_3=1679615$.
\end{theorem}

\begin{theorem}[Exactly $\{7,47\}$, and why]\label{thm:747}
For $n_j$ as above the associated candidate is $p_j=(4n_j+1)/3=2^{2^j+2}\,3^{2^j-1}-1$, and
$p_j\equiv7\pmod 8$ for all $j$. One has $p_0=7$ and $p_1=47$ prime, while
$p_2=1727=11\cdot157$ and $p_3=2239487=23\cdot 97369$. Moreover, for every prime $\ell\neq3$,
\[
  \ell\mid p \iff N\equiv-4^{-1}\!\!\pmod{\ell},
\]
and the two composite cases land exactly on such a killer residue:
$1295\equiv 8\equiv-4^{-1}\pmod{11}$ and $1679615\equiv17\equiv-4^{-1}\pmod{23}$.
\end{theorem}

\begin{corollary}[Unconditional bound]\label{cor:bound}
No exceptional prime exists below $1.33\times10^{14}$.
\end{corollary}

\noindent This improves Steinerberger's $10^{10}$ by a factor exceeding $13{,}000$, and is
obtained purely from Theorem~\ref{thm:bridge} together with Hercher's $n\ge10^{14}$.

\begin{theorem}[$\omega$-ladder]\label{thm:omega}
Every exceptional prime other than $7$ and $47$ arises from a solution $N$ with
$\omega(N)\ge6$; this is machine-checked in Lean~4 with no custom axiom. Certified exhaustive
enumeration of the $\omega=5,6,7$ strata (all empty; $272{,}676$ terminal nodes at $\omega=7$)
gives $\omega(N)\ge8$, beyond the published bound $\omega\ge7$ of \cite{Herch}.
\end{theorem}

\begin{theorem}[Congruence methods cannot close it]\label{thm:nogo}
For every prime $\ell\ge5$, every residue class of $N$ modulo $\ell$ compatible with the
equation is attained. Consequently no finite covering system of congruences forces $p$
composite, and the problem is not resolvable by covering or elementary congruence arguments
alone.
\end{theorem}

\subsection*{Formalization}
Theorems~\ref{thm:bridge}--\ref{thm:747} and the $\omega\le4$ classification are formalized in
Lean~4 against Mathlib. The cascade lemma (Theorem~\ref{thm:cascade}) has axiom footprint
exactly $\{\texttt{propext},\ \texttt{Classical.choice},\ \texttt{Quot.sound}\}$ --- no
\texttt{native\_decide}, no project-local axiom. See \S\ref{sec:lean}.

\section{The bridge}\label{sec:bridge}

\begin{proof}[Proof of Theorem~\ref{thm:bridge}]
From $N=(3p-1)/4$ we get $4N=3p-1$, hence $p=(4N+1)/3$. Then
\[
  \tot(N)=\frac{p+1}{2}
  =\frac{\frac{4N+1}{3}+1}{2}
  =\frac{4N+4}{6}
  =\frac{2N+2}{3},
\]
which is precisely $3\tot(N)=2N+2$, i.e.\ $\tot(N)=\tfrac23(N+1)$. Every step is reversible.
\end{proof}

\begin{proposition}[Master identity]\label{prop:master}
Every solution of $3\tot(N)=2N+2$ satisfies $p=(4N+1)/3=2\tot(N)-1$.
\end{proposition}

\begin{proof}
$3\tot(N)=2N+2$ gives $2N=3\tot(N)-2$, so $4N=6\tot(N)-4$ and
$4N+1=6\tot(N)-3=3\big(2\tot(N)-1\big)$. Divide by $3$.
\end{proof}

\begin{proof}[Proof of Theorem~\ref{thm:free}]
Let $N$ be odd with $N\equiv2\pmod 3$, so $3\mid 4N+1$ and $p=(4N+1)/3\in\Z$. Since
$3\cdot3=9\equiv1\pmod 8$, we have $3^{-1}\equiv3\pmod8$, whence
\[
  p\equiv 3(4N+1)=12N+3\equiv 4N+3 \pmod 8 .
\]
As $N$ is odd, $4N\equiv4\pmod 8$, so $p\equiv 7\pmod 8$.
\end{proof}

\begin{remark}
Theorem~\ref{thm:free} is worth emphasizing because the congruence $p\equiv7\pmod8$ appears in
the statement of \eqref{eq:exceptional} as though it were a constraint to be checked. It is
not: it is implied by the shape of $p$. Every solution $N$ of the totient equation is odd and
$\equiv2\pmod3$ (Proposition~\ref{prop:struct}), so the sole remaining question is whether
$p=2\tot(N)-1$ is prime.
\end{remark}

\begin{proposition}[Structure; cf.\ \cite{Herch}]\label{prop:struct}
Every solution $n$ of $3\tot(n)=2n+2$ is odd, satisfies $3\nmid n$, is squarefree, and is
\emph{chain-free}: no prime divisor $r$ of $n$ divides $s-1$ for a prime divisor $s$ of $n$.
Moreover every prime divisor of $n$ is $\ge5$.
\end{proposition}

\begin{proof}[Proof of chain-freeness]
Let $r,s\mid n$ be primes. From $s\mid n$ we get $\tot(s)=s-1\mid\tot(n)$. If $r\mid s-1$ then
$r\mid\tot(n)$, so $r\mid 3\tot(n)=2n+2$. But also $r\mid n$, hence $r\mid 2n$, and
subtracting, $r\mid2$. Thus $r=2$, contradicting oddness of $n$.
\end{proof}

\begin{proposition}[Sharper congruence]\label{prop:mod9}
Every solution satisfies $N\equiv5$ or $8\pmod 9$; in particular $3\nmid p$.
\end{proposition}

\begin{proposition}[$2$-adic structure]\label{prop:2adic}
$p\equiv-1 \pmod{2^{\omega(N)+1}}$, and $v_2(p-1)=1$ exactly. In the open regime
$\omega(N)\ge7$ this forces $p\equiv-1\pmod{256}$.
\end{proposition}

\section{The cascade}\label{sec:cascade}

\begin{lemma}[Telescoping]\label{lem:tele}
For any $x$ and $j\ge0$, $\ \prod_{k=0}^{j-1}\big(x^{2^k}+1\big)=\dfrac{x^{2^j}-1}{x-1}$.
In particular
\[
  6^{2^j}-1 \;=\; 5\cdot\prod_{k=0}^{j-1}\big(6^{2^k}+1\big),
\]
with the first factors $6+1=7$, $6^2+1=37$, $6^4+1=1297$.
\end{lemma}

\begin{proof}[Proof of Theorem~\ref{thm:cascade}]
Suppose $5$ and $6^{2^k}+1$ for $k<j$ are all prime. They are pairwise distinct, and each is
coprime to the others (consecutive Fermat-type factors differ by $2$ and are odd), so $n_j$ is
squarefree with
\[
  \tot(n_j)=4\prod_{k=0}^{j-1}6^{2^k}
  =4\cdot 6^{\,2^j-1}
  =\tfrac{4}{6}\cdot 6^{2^j}
  =\tfrac23\big(n_j+1\big),
\]
using $\sum_{k<j}2^k=2^j-1$. Hence $n_j$ solves the equation.

Conversely suppose $6^{2^m}+1$ is composite for some $m<j$, and take the least such. Write
$6^{2^j}-1=\big(6^{2^{j-1}}-1\big)\big(6^{2^{j-1}}+1\big)$; the two factors are odd and differ
by $2$, hence coprime, so $\tot$ is multiplicative across them. By induction the ``$-1$'' half
carries the exact value, while the ``$+1$'' half is divisible by the composite $6^{2^m}+1$ and
therefore has $\tot\big(6^{2^{j-1}}+1\big)<6^{2^{j-1}}$. Thus
$\tot(n_j)<\tfrac23(n_j+1)$ and $n_j$ is not a solution.

Finally $6^8+1=1679617=17\cdot98801$ is composite, so the cascade terminates at $j=3$.
\end{proof}

\begin{remark}
Base-$6$ generalized Fermat primes $6^{2^k}+1$ run out at $k=2$, exactly as the classical
Fermat primes run out at $F_4$. The finiteness of the known solution set is therefore not a
coincidence of small numbers but a structural fact about this cascade.
\end{remark}

\section{The primality sieve}\label{sec:sieve}

\begin{proof}[Proof of the sieve in Theorem~\ref{thm:747}]
Let $\ell\neq3$ be prime. Since $3p=4N+1$ and $\ell\nmid3$,
\[
  \ell\mid p \iff \ell\mid 4N+1 \iff 4N\equiv-1 \iff N\equiv-4^{-1} \pmod \ell . \qedhere
\]
\end{proof}

\noindent The two composite cascade candidates are explained exactly by this criterion:
\[
  1295\equiv 8 \pmod{11},\quad 4^{-1}\equiv3,\ -4^{-1}\equiv 8 \pmod{11}
  \;\Longrightarrow\; 11\mid p_2=1727,
\]
\[
  1679615\equiv 17 \pmod{23},\quad 4^{-1}\equiv6,\ -4^{-1}\equiv 17\pmod{23}
  \;\Longrightarrow\; 23\mid p_3=2239487 .
\]
In fact the criterion separates the four cascade members \emph{exactly}. Searching killer
residues $\ell\le\sqrt{p_j}$:
\[
\begin{array}{llll}
  N=5, & p=7, & \text{killers }\ell\le\sqrt p:\ \text{none} & \Rightarrow p\text{ prime},\\
  N=35, & p=47, & \text{none} & \Rightarrow p\text{ prime},\\
  N=1295, & p=1727, & \ell=11 & \Rightarrow p\text{ composite},\\
  N=1679615, & p=2239487, & \ell=23 & \Rightarrow p\text{ composite}.
\end{array}
\]
The two surviving cases avoid every killer residue below $\sqrt p$, and each failing case is
killed by exactly one. This is \emph{why} the cascade yields precisely $\{7,47\}$, rather than
a numerical accident.

\begin{proposition}[Intrinsic characterization]\label{prop:intrinsic}
A solution $n$ lies in the cascade if and only if $n+1$ is $\{2,3\}$-smooth, if and only if
$\tot(n)$ is $\{2,3\}$-smooth. The cascade primes $5,7,37,1297$ are exactly the primes $r$ for
which $r-1$ is $\{2,3\}$-smooth.
\end{proposition}

\section{The multiplicative-balance cap theorem}\label{sec:cap}

The $\omega$-ladder rests on a single inequality, which we isolate because it is independent of
number theory and drives the entire exhaustive search.

\begin{theorem}[Cap theorem]\label{thm:cap}
Let $A,B,c,r\in\N$ with $c\ge1$, and let $x_1,\dots,x_n$ be natural numbers with $x_i\ge c$ for
all $i$. If
\[
  B\,c^{\,n}+r \;<\; A\,(c-1)^{\,n},
\]
then
\[
  A\prod_{i=1}^{n}(x_i-1)\;=\;B\prod_{i=1}^{n}x_i+r
\]
is impossible. Equivalently: any solution of the balance must have some part $x_i<c$.
\end{theorem}

\begin{proof}
The map $x\mapsto (x-1)/x$ is increasing, in cross-multiplied form $(c-1)x\le(x-1)c$ whenever
$c\le x$. Taking the product over $i$,
\[
  (c-1)^n\prod_i x_i \;\le\; \Big(\prod_i (x_i-1)\Big)c^{\,n},
  \qquad\text{and}\qquad c^{\,n}\le\prod_i x_i .
\]
Write $Q=\prod x_i$, $P=\prod(x_i-1)$. Multiplying the hypothesis by $Q>0$ and the first
display by $A$, then substituting the balance $AP=BQ+r$,
\[
  (Bc^{\,n}+r)\,Q \;<\; A(c-1)^nQ \;\le\; A P\,c^{\,n}=(BQ+r)c^{\,n}.
\]
Expanding, the terms $Bc^nQ$ and $BQc^n$ agree and cancel, leaving $rQ<rc^{\,n}$, hence
$Q<c^{\,n}$ --- contradicting $c^{\,n}\le Q$.
\end{proof}

\begin{remark}
Theorem~\ref{thm:cap} is stated over a list of arbitrary length and with a free remainder $r$;
it is formalized in that generality (\S\ref{sec:lean}). It contains no prime, totient, or
divisibility hypothesis. Its role is to prune: a single numeric check kills an infinite branch
of the search, which is what converts an enumeration into a proof of non-existence. Any system
carrying a multiplicative balance against a monotone per-part discount inherits it.
\end{remark}

\section{The \texorpdfstring{$\omega$}{omega}-ladder}\label{sec:omega}

Combining Proposition~\ref{prop:struct} with Theorem~\ref{thm:cap} yields a finite decision
tree at each stratum $\omega(N)=w$: the cap theorem bounds the smallest remaining prime, the
chain-freeness condition prunes candidates, and two-variable terminals are resolved by an exact
divisor identity. Levels $w\le4$ are classified completely and unconditionally; the solutions
are exactly $5,\,35,\,1295,\,1679615$. Along the cascade branch the terminal eliminations take
the strikingly rigid form
\[
  (p-3)(q-3)=8,\qquad (q-6)(r-6)=31,\qquad (r-36)(s-36)=1261=13\cdot 97,
\]
realized by $(5,7)$, $(7,37)$ and $(37,1297)$ respectively --- in each case the left factor is
forced to $1$ (or $2$) and the right factor is pinned to a divisor of a fixed integer, which is
what closes the level. Other branches carry their own eliminations of the same shape.

\begin{theorem}\label{thm:w6}
Any exceptional prime other than $7$ and $47$ arises from a solution with $\omega(N)\ge6$.
This is verified in the Lean kernel.
\end{theorem}

\noindent The strata $\omega=5,6,7$ are certified empty by exhaustive enumeration
($22$, $410$, and $272{,}676$ terminal nodes respectively). The $\omega=7$ computation was
independently re-run from scratch with separate code and no chain filtering --- strictly more
permissive --- reproducing the terminal census exactly and finding nothing; every equation-
satisfying integer pair in the $\omega=7$ universe fails on compositeness alone, with explicit
witnesses. This yields $\omega(N)\ge8$ for any further solution, past the published $\omega\ge7$
of \cite{Herch}. We label this result \emph{computational}: it is exhaustive and cross-checked,
but the $\omega=6,7$ strata are not yet reproduced inside the Lean kernel.

\section{A no-go for congruence methods}\label{sec:nogo}

\begin{proof}[Proof sketch of Theorem~\ref{thm:nogo}]
Fix a prime $\ell\ge5$. The equation $3\tot(N)=2N+2$ places no obstruction on the residue of
$N$ modulo $\ell$ beyond those already imposed by squarefreeness and chain-freeness, and each
admissible class is realized. Hence for no finite set of moduli does every residue class force
$N\equiv-4^{-1}$ modulo some member, i.e.\ no covering system forces $p$ composite.
\end{proof}

\noindent The consequence is methodological: the difficulty is of Fermat-prime type --- an
existence question about primes in a sparse family --- and cannot be dissolved by local
information. A Mertens-product heuristic suggests that ``almost all'' further solutions would
yield composite $p$; we stress that this is a heuristic and not a proof.

\section{Formal verification}\label{sec:lean}

All statements below are checked against Lean~4 with Mathlib.

\begin{center}
\begin{tabular}{@{}lll@{}}
\hline
Statement & Lean name & Axiom footprint\\
\hline
Reduction (Thm.~\ref{thm:bridge}) & \texttt{reduction} & $\{$propext, choice, Quot.sound$\}$\\
Master identity (Prop.~\ref{prop:master}) & \texttt{master\_identity} & $\{$propext, choice, Quot.sound$\}$\\
Free residue clause (Thm.~\ref{thm:free}) & \texttt{mod8\_free} & $\{$propext, choice, Quot.sound$\}$\\
Cascade lemma (Thm.~\ref{thm:cascade}) & \texttt{cascade\_lemma} & $\{$propext, choice, Quot.sound$\}$\\
Structure (Prop.~\ref{prop:struct}) & \texttt{solution\_squarefree} etc. & $\{$propext, choice, Quot.sound$\}$\\
$\omega\le4$ classification & \texttt{solution\_omega\_le\_four\_classified} & $\{$propext, choice, Quot.sound$\}$\\
$\omega\ge6$ (Thm.~\ref{thm:w6}) & \texttt{exceptional\_high\_omega\_six} & $+$ \texttt{native\_decide}\\
Cap theorem (Thm.~\ref{thm:cap}) & \texttt{cap\_kill\_balance} & $\{$propext, choice, Quot.sound$\}$\\
\hline
\end{tabular}
\end{center}

\noindent The cascade lemma and the cap theorem carry no \texttt{native\_decide} and no
project-local axiom. The $\omega\ge6$ result inherits disclosed \texttt{native\_decide}
compiler-trust axioms from the finite scans in its kill-tree.

\section{What remains open}

\begin{conjecture}[Steinerberger]
The only solutions of $\tot(n)=\tfrac23(n+1)$ are $5$, $35$, $1295$, $1679615$.
\end{conjecture}

The $r=2$ case of Erd\H{o}s--Graham \#411 is complete if and only if no \emph{non-cascade}
solution of $\tot(n)=\tfrac23(n+1)$ --- necessarily squarefree, chain-free, with at least eight
distinct prime factors all $\ge5$, and exceeding $10^{14}$ --- yields a prime $(4n+1)/3$. By
Theorem~\ref{thm:nogo} this cannot be settled by congruence methods. We record the precise
remaining statement:

\begin{quote}
\emph{Every solution $N$ of $3\tot(N)=2N+2$ with $N>1679615$ has $(4N+1)/3$ composite.}
\end{quote}

A proof of that statement closes the problem; nothing weaker is required, and nothing stronger
is needed.

\section*{Acknowledgements}
The computations were carried out with a deterministic, zero-cost verification pipeline; all
certificates, kill-trees, and Lean sources are available from the author.

\begin{thebibliography}{9}

\bibitem{Stein} S.~Steinerberger, \emph{On an iterated arithmetic function problem of
Erd\H{o}s and Graham}, arXiv:2504.08023 (2025).

\bibitem{Herch} C.~Hercher, \emph{On positive integers $n$ with $\varphi(n)=\tfrac23(n+1)$},
arXiv:2504.19915 (2025).

\bibitem{EG} P.~Erd\H{o}s and R.~L.~Graham, \emph{Old and New Problems and Results in
Combinatorial Number Theory}, Monographies de L'Enseignement Math\'ematique 28, Gen\`eve, 1980.

\bibitem{Mathlib} The mathlib Community, \emph{The Lean mathematical library},
CPP 2020, 367--381.

\end{thebibliography}

\end{document}
